\begin{table}[htbp]
\centering
\caption{Resumo da árvore de decisão multinomial, recorte geral, avaliação in\_sample, profundidade máxima 19.}
\label{tab:tabela_compacta_treeClassification_19_profundidade_ternario_recorte_geral_in_sample}
\small
\begin{tabular}{lllllll}
\toprule
\rowcolor[gray]{0.85}
Especificação & N & ROC-AUC & Renda & Desempenho & Interação & Variáveis \\
\midrule
Modelo 1 & 1.925.666 & 0,7676 & 0,1128 & 0,1008 & 0,7863 & 3 \\
\rowcolor[gray]{0.95}
Modelo 2 & 1.925.666 & 0,8524 & 0,0277 & 0,0800 & 0,3335 & 6 \\
Modelo 3 & 1.925.666 & 0,9053 & 0,0218 & 0,0594 & 0,2678 & 10 \\
\rowcolor[gray]{0.95}
Modelo 4 & 1.925.666 & 0,9139 & 0,0172 & 0,0686 & 0,2555 & 14 \\
Modelo 5 & 1.925.666 & 0,9216 & 0,0161 & 0,0671 & 0,2485 & 20 \\
\bottomrule
\end{tabular}
\par\footnotesize{Notas: Renda, desempenho e interação indicam importâncias normalizadas das variáveis na árvore. A profundidade máxima parametrizada é 19.}
\par\footnotesize{Fonte: elaboração própria (2026).}
\end{table}
